% =====================================================================
% CHƯƠNG 2 - CƠ SỞ LÝ THUYẾT
% =====================================================================
\chapter{CƠ SỞ LÝ THUYẾT}

Chương này tổ chức cơ sở lý thuyết theo đúng pipeline của đề tài: từ đặc trưng
Edge AI, mô hình YOLO, các kỹ thuật nén và lượng tử hóa, TensorRT và Jetson đến
quản lý model deployment, Edge Runtime, nền tảng phân tán, RAG, AI Agent và
Domain-Driven Design. Cách sắp xếp này tạo cầu nối trực tiếp từ kiến thức nền
tảng sang nghiên cứu tối ưu ở Chương 3 và thiết kế nền tảng ở Chương 4.

\section{Edge AI}

\subsection{Khái niệm Edge AI}

Edge AI là cách triển khai tác vụ trí tuệ nhân tạo trên thiết bị đặt gần nguồn sinh
dữ liệu, thay vì gửi toàn bộ dữ liệu thô về đám mây. Thiết bị có thể là gateway,
máy tính nhúng, camera thông minh hoặc máy công nghiệp. Lợi ích chính gồm giảm
độ trễ, giảm băng thông, tăng khả năng hoạt động ngoại tuyến và hạn chế dữ liệu
nhạy cảm rời khỏi hiện trường \cite{shi2016,satya2017}.

Edge AI không loại bỏ vai trò của Cloud. Cloud vẫn phù hợp cho huấn luyện, lưu
trữ artifact, quản lý đội thiết bị và phân tích dài hạn; Edge chịu trách nhiệm cho
vòng lặp phản hồi nhanh. Do đó, giải pháp thực tế thường là một hệ thống lai trong
đó Cloud và Edge có nhiệm vụ khác nhau nhưng chia sẻ cùng vòng đời model.

\subsection{Kiến trúc Cloud--Edge}

Kiến trúc Cloud--Edge gồm ba lớp chính: Cloud Control Plane, kênh truyền thông
và Edge Plane. Control Plane lưu desired state, model version, chính sách và lịch
sử vận hành. Edge Plane chứa Edge Agent và AI Runtime, duy trì actual state và
tiếp tục suy luận khi mất kết nối. Kênh truyền thông mang lệnh, sự kiện và
telemetry; các tệp model dung lượng lớn có thể được phân phối qua object storage.

Kiến trúc cần chấp nhận kết nối không ổn định. Lệnh phải có định danh duy nhất,
thời hạn hiệu lực và khả năng thực thi lặp an toàn. Edge Agent phải lưu trạng thái
cục bộ, đồng bộ lại khi có mạng và không được dừng runtime đang khỏe chỉ vì
Control Plane tạm thời không truy cập được.

\subsection{Đặc điểm của Edge Device}

Edge Device có tài nguyên nhỏ hơn máy chủ, thường dùng bộ nhớ chia sẻ giữa CPU
và GPU, giới hạn điện năng và tản nhiệt. Phần mềm phụ thuộc chặt vào phiên bản
hệ điều hành, CUDA, cuDNN, TensorRT và driver. Cùng một model artifact có thể
không chạy được trên hai thiết bị nếu engine được tạo cho kiến trúc GPU hoặc
phiên bản runtime khác nhau.

Thiết bị biên cũng nằm ngoài trung tâm dữ liệu nên dễ gặp mất nguồn, mất mạng,
camera lỗi hoặc thay đổi môi trường. Vì vậy, khả năng quan sát, tự khởi động lại,
rollback và xác thực artifact là một phần của yêu cầu hệ thống chứ không phải chức
năng bổ sung.

\subsection{Các vấn đề khi triển khai AI trên Edge}

Năm nhóm vấn đề chính là năng lực tính toán, bộ nhớ, latency, power và thermal.
Chúng tạo thành bài toán tối ưu đa mục tiêu: giảm số phép tính nhưng không làm
accuracy giảm quá mức; giảm precision nhưng vẫn duy trì tính ổn định số; tăng FPS
nhưng không gây quá nhiệt hoặc vượt ngân sách điện.

Ngoài model inference, hệ thống còn phải tính chi phí camera capture, chuyển đổi
màu, resize, chuẩn hóa, non-maximum suppression, truyền kết quả và telemetry.
Do đó, benchmark chỉ đo thời gian TensorRT engine là cần thiết nhưng chưa đủ để
kết luận về khả năng đáp ứng thời gian thực của toàn hệ thống.
\section{Object Detection và YOLO}

\subsection{Bài toán Object Detection}

Object Detection xác định đồng thời lớp và vị trí của các đối tượng trong ảnh.
Đầu ra thường là tập các bounding box $b_i=(x_i,y_i,w_i,h_i)$, nhãn lớp và độ
tin cậy. Khác với phân loại ảnh, một ảnh có thể chứa nhiều đối tượng và mô hình
phải giải quyết cả bài toán định vị lẫn phân loại.

Độ chồng lấp giữa hộp dự đoán $A$ và hộp chuẩn $B$ được đo bằng Intersection
over Union:
\begin{equation}
  IoU(A,B)=\frac{|A\cap B|}{|A\cup B|}.
\end{equation}
Một dự đoán được xem là đúng khi lớp phù hợp và IoU vượt ngưỡng đánh giá.

\subsection{Kiến trúc YOLO}

YOLO là họ mô hình phát hiện đối tượng một giai đoạn, thực hiện trích xuất đặc
trưng và dự đoán trong một đồ thị thống nhất \cite{redmon2016,redmon2017}.
Kiến trúc tổng quát gồm backbone để trích xuất đặc trưng, neck để tổng hợp đặc
trưng đa tỉ lệ và detection head để dự đoán vị trí, objectness và lớp.

Các biến thể YOLO khác nhau ở block tích chập, cơ chế kết hợp đặc trưng, hàm
loss và cách biểu diễn detection head. Khi tối ưu, cần bảo toàn những tầng ảnh
hưởng mạnh đến đầu ra, đồng thời ưu tiên cắt tỉa cấu trúc mà TensorRT có thể khai
thác thành giảm thời gian thực thi thực tế.

\subsection{Quy trình huấn luyện YOLO}

Quy trình huấn luyện gồm chuẩn hóa nhãn, chia train--validation--test, tăng cường
dữ liệu, khởi tạo trọng số, tối ưu loss và chọn checkpoint. Mọi thí nghiệm phải
lưu seed, phiên bản mã nguồn, cấu hình optimizer, learning rate, batch size, image
size và số epoch để có thể tái lập.

Tập test chỉ được dùng cho đánh giá cuối. Các quyết định về pruning ratio, lịch
fine-tuning và ngưỡng dừng sớm phải dựa trên tập validation nhằm tránh điều chỉnh
gián tiếp theo tập test.

\subsection{Các chỉ số đánh giá}

Với $TP$, $FP$ và $FN$ lần lượt là số dự đoán đúng, dương tính giả và âm tính
giả, Precision và Recall được xác định bởi:
\begin{equation}
  Precision=\frac{TP}{TP+FP}, \qquad
  Recall=\frac{TP}{TP+FN}.
\end{equation}

Average Precision là diện tích dưới đường Precision--Recall của một lớp. mAP
là trung bình AP trên các lớp. $mAP@50$ dùng ngưỡng IoU 0,5; $mAP@50{:}95$
lấy trung bình trên các ngưỡng từ 0,5 đến 0,95 với bước 0,05. Đồ án báo cáo đồng
thời Precision, Recall, $mAP@50$ và $mAP@50{:}95$ để tránh kết luận dựa trên
một chỉ số đơn lẻ.
\section{Model Compression}

\subsection{Tổng quan Model Compression}

Model Compression là nhóm kỹ thuật giảm chi phí lưu trữ hoặc suy luận của mạng
nơ-ron. Các hướng phổ biến gồm pruning, quantization, knowledge distillation và
thiết kế kiến trúc gọn. Đồ án tập trung vào pruning và quantization vì hai kỹ thuật
này có thể áp dụng trực tiếp lên baseline và kết hợp với TensorRT.

Hiệu quả nén phải được đánh giá trên cả kích thước artifact, số phép tính lý
thuyết và latency thực tế. Một mô hình có nhiều trọng số bằng không chưa chắc
chạy nhanh hơn nếu backend vẫn thực hiện phép tính dense.

\subsection{Model Pruning}

Pruning loại bỏ trọng số, kênh hoặc cấu trúc được xem là ít quan trọng. Tiêu chí
có thể dựa trên độ lớn trọng số, chuẩn của filter, độ nhạy loss hoặc thông tin
gradient. Quy trình cơ bản gồm huấn luyện baseline, tính điểm quan trọng, cắt tỉa,
fine-tuning và đánh giá lại \cite{han2016,molchanov2017}.

\subsection{Structured Pruning}

Structured pruning loại bỏ đơn vị có cấu trúc như channel, filter hoặc block.
Kết quả là tensor có kích thước nhỏ hơn và có thể được thư viện dense hiện có
khai thác mà không cần kernel thưa chuyên dụng. Đây là lựa chọn phù hợp khi mục
tiêu cuối là TensorRT trên Jetson.

Nhược điểm là mỗi quyết định loại bỏ channel ảnh hưởng đến các tầng lân cận và
có thể làm accuracy giảm nhanh. Việc cắt tỉa cần tôn trọng quan hệ residual,
concatenation và detection head trong kiến trúc YOLO.

\subsection{Unstructured Pruning}

Unstructured pruning đặt các trọng số riêng lẻ về không và có thể đạt sparsity
cao với mức suy giảm accuracy thấp hơn. Tuy nhiên, kích thước tensor không đổi;
latency chỉ giảm rõ khi phần cứng và runtime hỗ trợ sparse kernel tương ứng.

Trong đồ án, unstructured pruning có thể được dùng làm đối chứng về sparsity,
nhưng cấu hình triển khai ưu tiên structured pruning nếu mục tiêu là giảm latency
thực tế trên Jetson.

\subsection{Pruning Ratio}

Pruning ratio biểu diễn tỷ lệ phần tử hoặc cấu trúc bị loại bỏ:
\begin{equation}
  r=\frac{N_{before}-N_{after}}{N_{before}}\times 100\%.
\end{equation}
Các mức 20\%, 30\% và 40\% được dùng làm điểm thực nghiệm ban đầu. Tỷ lệ này
không mặc định áp dụng đồng đều cho mọi tầng; các tầng đầu, tầng đầu ra và nhánh
nhạy cảm có thể được bảo vệ hoặc đặt ngưỡng riêng.

\subsection{Fine-tuning sau Pruning}

Sau pruning, phân bố đặc trưng bị thay đổi nên mô hình cần được fine-tune với
learning rate nhỏ hơn baseline. Fine-tuning phục hồi accuracy trong khi duy trì
cấu trúc đã cắt. Checkpoint tốt nhất được chọn theo validation mAP và phải được
đánh giá lại trên tập test trước khi chuyển sang QAT hoặc TensorRT.
\section{Model Quantization}

\subsection{Khái niệm Quantization}

Quantization ánh xạ giá trị số thực có độ chính xác cao sang biểu diễn có số bit
thấp hơn. Với lượng tử tuyến tính, một giá trị $x$ được ánh xạ thành số nguyên
$q$ thông qua scale $s$ và zero-point $z$:
\begin{equation}
  q=\operatorname{clip}\left(\operatorname{round}\left(\frac{x}{s}\right)+z\right).
\end{equation}
Việc giảm số bit làm giảm dung lượng và có thể tăng thông lượng, nhưng gây sai số
lượng tử cần được kiểm soát.

\subsection{FP32, FP16 và INT8}

FP32 là biểu diễn phổ biến khi huấn luyện và được dùng làm chuẩn chất lượng.
FP16 giảm một nửa dung lượng mỗi phần tử, phù hợp với Tensor Core và thường giữ
accuracy gần FP32. INT8 dùng số nguyên 8 bit, có tiềm năng giảm bộ nhớ và tăng
tốc mạnh hơn nhưng nhạy với dải giá trị activation và yêu cầu calibration hoặc
QAT.

\subsection{Post-Training Quantization}

Post-Training Quantization thực hiện lượng tử sau khi mô hình đã huấn luyện. PTQ
không yêu cầu huấn luyện lại hoặc chỉ cần calibration dataset để ước lượng dải
activation. Ưu điểm là quy trình nhanh; hạn chế là accuracy có thể giảm với tầng
nhạy cảm, dữ liệu lệch hoặc mô hình nhỏ.

\subsection{Quantization-Aware Training}

QAT mô phỏng hiệu ứng lượng tử trong quá trình huấn luyện để trọng số thích nghi
với sai số làm tròn và clipping \cite{jacob2018}. Forward pass sử dụng các nút
fake quantization, còn gradient được xấp xỉ để cập nhật trọng số thực. QAT tốn
thêm thời gian huấn luyện nhưng thường duy trì accuracy INT8 tốt hơn PTQ.

\subsection{Fake Quantization}

Fake quantization lượng tử và giải lượng tử tensor trong forward pass nhưng vẫn
lưu trọng số ở dạng số thực để tối ưu. Các observer thu thập thống kê min--max
hoặc histogram nhằm xác định scale. Khi export, các thông tin này được chuyển
thành toán tử lượng tử hoặc metadata phục vụ TensorRT.

\subsection{So sánh PTQ và QAT}

\begin{table}[H]
\centering
\caption{So sánh PTQ và QAT}
\label{tab:ptq-qat}
\begin{tabularx}{\textwidth}{L{3.2cm}YY}
\toprule
\textbf{Tiêu chí} & \textbf{PTQ} & \textbf{QAT}\\
\midrule
Dữ liệu cần thiết & Tập calibration đại diện & Dữ liệu huấn luyện/validation\\
Chi phí thực hiện & Thấp, không huấn luyện lại & Cao hơn do fine-tuning có fake quant\\
Khả năng giữ accuracy & Phụ thuộc độ nhạy mô hình & Thường tốt hơn với INT8\\
Độ phức tạp pipeline & Đơn giản & Cần cấu hình và export tương thích\\
Vai trò trong đồ án & Baseline INT8 tham chiếu & Phương pháp tối ưu chính\\
\bottomrule
\end{tabularx}
\end{table}
\section{TensorRT}

\subsection{Kiến trúc TensorRT}

TensorRT là bộ SDK tối ưu và thực thi suy luận cho GPU NVIDIA. Builder tiếp nhận
đồ thị mạng, áp dụng tối ưu như fusion, lựa chọn kernel và precision, sau đó tạo
serialized engine. Runtime nạp engine, cấp phát tensor và thực thi qua execution
context \cite{tensorrt}.

Engine phụ thuộc vào kiến trúc GPU, phiên bản TensorRT, plugin và shape profile.
Vì vậy, manifest của model artifact phải ghi rõ môi trường tạo engine và điều kiện
tương thích.

\subsection{ONNX}

ONNX là định dạng trung gian giúp trao đổi mô hình giữa framework huấn luyện và
runtime suy luận \cite{onnx}. Pipeline của đồ án export checkpoint sang ONNX,
kiểm tra graph, so sánh đầu ra với framework gốc rồi mới xây dựng TensorRT
engine. Các toán tử không được hỗ trợ phải được thay thế, phân rã hoặc hiện thực
bằng plugin.

\subsection{TensorRT Engine}

TensorRT engine là artifact đã được tối ưu cho một tập cấu hình cụ thể. Quá trình
build có thể lựa chọn tactic khác nhau theo workspace và phần cứng. Engine cần
đi kèm input shape, data type, tên binding, preprocessing, postprocessing và hàm
băm để Edge Runtime tải đúng cách.

\subsection{FP16 Inference}

FP16 inference cho phép TensorRT dùng kernel bán chính xác ở các tầng được hỗ
trợ và quay về precision khác khi cần. Đây thường là cấu hình triển khai đầu tiên
vì quy trình chuyển đổi đơn giản hơn INT8 và mức suy giảm accuracy nhỏ. Kết quả
vẫn phải được so sánh với FP32 trên cùng tập test.

\subsection{INT8 Inference}

INT8 inference yêu cầu scale activation phù hợp. Với PTQ, TensorRT dùng
calibrator; với QAT, thông tin lượng tử được mang theo từ graph nếu pipeline
export hỗ trợ. Một số tầng có thể giữ FP16 hoặc FP32 để bảo toàn accuracy, tạo ra
engine mixed precision.

\subsection{TensorRT trên NVIDIA Jetson}

Trên Jetson, TensorRT được phân phối cùng hệ sinh thái JetPack và sử dụng CUDA,
cuDNN cùng driver tương ứng. Việc build engine trực tiếp trên thiết bị đích giúp
giảm rủi ro không tương thích. Đồ án lưu riêng engine cho Jetson Nano và Jetson
Orin Nano, thay vì giả định một engine có thể dùng chung cho mọi GPU.
\section{NVIDIA Jetson}

\subsection{Jetson Nano}

Jetson Nano đại diện cho nhóm thiết bị Edge AI có tài nguyên hạn chế. Thiết bị
phù hợp để đánh giá mức tối thiểu mà model có thể vận hành, đồng thời làm rõ ảnh
hưởng của bộ nhớ, power mode và thermal throttling. Phiên bản JetPack, CUDA,
TensorRT, dung lượng RAM và chế độ công suất phải được ghi theo đúng thiết bị
thực nghiệm, vì các biến thể phần cứng có thể khác nhau.

\subsection{Jetson Orin Nano}

Jetson Orin Nano đại diện cho thế hệ GPU nhúng mới hơn, có kiến trúc và năng lực
tăng tốc khác Jetson Nano. Thiết bị được dùng để đánh giá khả năng mở rộng hiệu
năng của cùng pipeline model, đồng thời kiểm tra tính di động của Model Registry,
Deployment và Edge Runtime giữa hai lớp node.

\subsection{So sánh tài nguyên phần cứng}

\begin{table}[H]
\centering
\caption{Thông tin phần cứng cần ghi nhận cho hai thiết bị Jetson}
\label{tab:jetson-hardware}
\begin{tabularx}{\textwidth}{L{4cm}YY}
\toprule
\textbf{Thành phần} & \textbf{Jetson Nano} & \textbf{Jetson Orin Nano}\\
\midrule
Kiến trúc GPU & \cbs{thông tin thiết bị} & \cbs{thông tin thiết bị}\\
Số nhân GPU/Tensor Core & \cbs{thông tin thiết bị} & \cbs{thông tin thiết bị}\\
RAM & \cbs{dung lượng} & \cbs{dung lượng}\\
JetPack & \cbs{phiên bản} & \cbs{phiên bản}\\
CUDA & \cbs{phiên bản} & \cbs{phiên bản}\\
TensorRT & \cbs{phiên bản} & \cbs{phiên bản}\\
Power mode & \cbs{chế độ đo} & \cbs{chế độ đo}\\
Tản nhiệt & \cbs{cấu hình} & \cbs{cấu hình}\\
\bottomrule
\end{tabularx}
\end{table}

Bảng~\ref{tab:jetson-hardware} phải được điền từ lệnh hệ thống và nhãn thiết bị
thực tế, không suy ra từ thông số quảng cáo của một phiên bản khác.

\subsection{CUDA và GPU Acceleration}

CUDA cung cấp mô hình lập trình và runtime để thực thi kernel trên GPU NVIDIA.
Trong pipeline, TensorRT sử dụng CUDA để lập lịch kernel và truyền dữ liệu. Tối
ưu end-to-end cần hạn chế sao chép CPU--GPU, tái sử dụng bộ đệm, dùng stream
phù hợp và đồng bộ chỉ tại những điểm cần thiết.

GPU acceleration không tự động bảo đảm latency thấp. Kích thước input, batch,
memory allocation, camera decode và hậu xử lý trên CPU có thể trở thành nút thắt.
Vì vậy, đồ án đo từng giai đoạn của pipeline và tổng thời gian đầu cuối.
\section{AI Model Deployment}

\subsection{Model Registry}

Model Registry là kho quản lý metadata và artifact của các phiên bản mô hình.
Mỗi bản ghi liên kết model name, version, framework, precision, input shape, chỉ
số đánh giá, vị trí artifact, checksum và ma trận tương thích thiết bị. Registry
không chỉ là thư mục lưu tệp mà là nguồn sự thật cho quyết định triển khai.

\subsection{Model Versioning}

Model versioning giúp phân biệt các checkpoint, cấu hình tối ưu và engine. Một
phiên bản đã phát hành phải bất biến; thay đổi nội dung yêu cầu tạo phiên bản mới.
Ngoài version dễ đọc, hệ thống lưu checksum để nhận diện chính xác bytes của
artifact và ngăn trường hợp hai tệp khác nhau dùng chung tên.

\subsection{Model Artifact}

Artifact triển khai có thể gồm TensorRT engine, ONNX dự phòng, label map,
manifest, cấu hình preprocessing--postprocessing và chữ ký hoặc checksum. Gói
phải tự mô tả đủ để Edge Agent kiểm tra tính tương thích trước khi dừng runtime
hiện tại.

\subsection{Deployment Strategy}

Chiến lược triển khai xác định thứ tự và phạm vi node nhận model. Các lựa chọn
gồm triển khai đồng loạt, rolling deployment và canary. Với đội thiết bị phân tán,
canary giúp kiểm tra một nhóm nhỏ trước khi mở rộng, đồng thời đặt ngưỡng dừng
dựa trên health check, FPS, error rate hoặc nhiệt độ.

\subsection{Rollback}

Rollback đưa node về phiên bản đã biết là ổn định khi download, validation,
activation hoặc health check thất bại. Edge Agent cần giữ ít nhất một artifact
trước đó, cập nhật con trỏ active theo cách nguyên tử và báo cáo nguyên nhân. Cơ
chế rollback phải được kiểm thử như một luồng bình thường, không chỉ dùng khi
sự cố đã xảy ra.
\section{Edge Runtime}

\subsection{AI Inference Pipeline}

Edge Runtime biến frame từ camera thành detection theo chuỗi trong
Hình~\ref{fig:inference-pipeline-theory}.

\begin{figure}[H]
\centering
\begin{tikzpicture}[node distance=0.75cm]
  \node[flowbox] (camera) {Camera};
  \node[flowbox, below=of camera] (pre) {Preprocess};
  \node[flowbox, below=of pre] (infer) {Inference};
  \node[flowbox, below=of infer] (post) {Postprocess};
  \node[flowbox, below=of post] (result) {Result};
  \draw[flowarrow] (camera) -- (pre);
  \draw[flowarrow] (pre) -- (infer);
  \draw[flowarrow] (infer) -- (post);
  \draw[flowarrow] (post) -- (result);
\end{tikzpicture}
\caption{Pipeline suy luận tại Edge}
\label{fig:inference-pipeline-theory}
\end{figure}

Preprocess phải khớp với huấn luyện về resize, letterbox, thứ tự kênh và chuẩn
hóa. Postprocess giải mã tensor, áp dụng confidence threshold và non-maximum
suppression. Sai khác ở hai bước này có thể làm giảm accuracy dù engine cho đầu
ra số học chính xác.

\subsection{Runtime Lifecycle}

Runtime lifecycle gồm các trạng thái created, loading, ready, running, degraded,
stopping, stopped và failed. Mọi chuyển trạng thái phải có điều kiện hợp lệ và
reason code. Edge Agent không coi tiến trình đang tồn tại là đủ; health check cần
xác nhận camera, engine và vòng xử lý frame đều hoạt động.

\subsection{GPU Resource Management}

Quản lý GPU bao gồm cấp phát bộ nhớ, số execution context, stream, batch size và
số runtime đồng thời. Hệ thống cần từ chối hoặc trì hoãn cấu hình vượt khả năng
node thay vì để thiết bị rơi vào out-of-memory. Model manifest và node capability
được dùng để kiểm tra trước khi activation.

\subsection{Runtime Monitoring}

Runtime phát sinh các chỉ số FPS, latency từng stage, queue depth, frame drop,
GPU usage, GPU memory, nhiệt độ, power và error count. Metrics cần gắn node,
runtime, model version và timestamp. Cửa sổ tổng hợp giúp phát hiện suy giảm kéo
dài thay vì tạo incident từ một mẫu nhiễu đơn lẻ.
\section{Distributed Edge Platform}

\subsection{Cloud Control Plane}

Cloud Control Plane duy trì trạng thái mong muốn và điều phối các domain nghiệp
vụ. Thành phần này cung cấp API cho Dashboard và Copilot, lưu lịch sử model,
deployment, command, incident và audit. Control Plane không tham gia trực tiếp
vào vòng lặp inference nên Edge vẫn hoạt động khi Cloud tạm thời gián đoạn.

\subsection{Edge Agent}

Edge Agent là đại diện quản lý trên mỗi node. Agent đăng ký danh tính, nhận lệnh,
đồng bộ desired state, tải artifact, điều khiển runtime và gửi actual state. Agent
phải giới hạn quyền hệ thống, xác thực nguồn lệnh và bảo đảm một command không
bị thực thi hai lần khi message broker gửi lại.

\subsection{Telemetry}

Telemetry gồm metric, log, event và state report. Metric phục vụ theo dõi xu
hướng; event ghi nhận thay đổi rời rạc; log hỗ trợ chẩn đoán chi tiết; state report
đối chiếu desired state với actual state. Khi mất mạng, dữ liệu quan trọng được
đệm cục bộ theo dung lượng giới hạn và gửi lại theo thứ tự phù hợp.

\subsection{Message Broker}

Message broker tách Control Plane khỏi kết nối trực tiếp đến từng node. MQTT phù
hợp với thiết bị biên nhờ mô hình publish--subscribe và mức QoS linh hoạt
\cite{mqtt5}. Topic phải được phân vùng theo tenant/node, áp dụng TLS và quyền
publish/subscribe tối thiểu.

\subsection{Monitoring}

Monitoring thu thập, lưu trữ và trực quan hóa telemetry; đồng thời đánh giá rule
để phát hiện bất thường vận hành. Một rule cần metric, cửa sổ thời gian, ngưỡng,
mức nghiêm trọng và điều kiện đóng. Kết quả phát hiện được chuyển thành Incident
thay vì chỉ tồn tại dưới dạng biểu đồ.
\section{RAG và AI Agent}

\subsection{Retrieval-Augmented Generation}

Retrieval-Augmented Generation (RAG) kết hợp truy xuất tài liệu với mô hình
ngôn ngữ. Câu hỏi được dùng để tìm các đoạn liên quan; nội dung truy xuất được
đưa vào ngữ cảnh để tạo câu trả lời có căn cứ \cite{lewis2020rag}. Trong đồ án,
Knowledge chứa runbook, hướng dẫn triển khai, mô tả lỗi và tài liệu hệ thống.

\subsection{Embedding Search}

Embedding search ánh xạ truy vấn và tài liệu thành vector rồi xếp hạng theo độ
tương đồng. Phương pháp này tìm được nội dung gần nghĩa dù từ khóa không trùng
khớp. Chất lượng phụ thuộc cách chia đoạn, mô hình embedding, metadata filter và
ngưỡng lấy kết quả.

\subsection{Hybrid Search}

Hybrid search kết hợp tìm kiếm từ khóa với vector search. Từ khóa phù hợp với mã
lỗi, model version và tên node; vector search phù hợp với câu hỏi diễn đạt tự
nhiên. Hai danh sách có thể được chuẩn hóa điểm, hợp nhất và rerank trước khi đưa
cho Agent.

\subsection{AI Agent}

AI Agent mở rộng RAG bằng khả năng lập kế hoạch và gọi công cụ. Agent có thể
truy vấn Fleet, Runtime, Deployment hoặc Incident thay vì chỉ trả lời từ tài liệu.
Mô hình ngôn ngữ không được truy cập trực tiếp cơ sở dữ liệu hay shell; mọi hành
động đi qua tool schema, authorization và audit.

\subsection{Tool Calling}

Tool calling là cơ chế mô hình chọn hàm và tạo tham số có cấu trúc. Hệ thống phải
kiểm tra schema, quyền người dùng, phạm vi node và policy trước khi thực thi. Các
tool thay đổi trạng thái dùng hai pha: tạo kế hoạch để người dùng xem xét, sau đó
phê duyệt bằng định danh kế hoạch và token chống thực thi lại.

\subsection{MCP}

Model Context Protocol (MCP) chuẩn hóa cách AI application khám phá và gọi
tools/resources do server cung cấp. Trong phạm vi đồ án, MCP được dùng làm biên
tích hợp cho Knowledge và các công cụ vận hành. MCP không thay thế cơ chế xác
thực hoặc phân quyền của domain; server vẫn phải xác minh danh tính và chính
sách cho từng lần gọi.
\section{Domain-Driven Design}

Domain-Driven Design (DDD) tổ chức phần mềm quanh ngôn ngữ và quy tắc nghiệp
vụ thay vì quanh bảng dữ liệu hoặc framework \cite{evans2003ddd}. Cách tiếp cận
này phù hợp với nền tảng có nhiều miền như Fleet, Model Registry, Deployment,
Incident, Knowledge, Copilot và Ops.

\subsection{Entity}

Entity có định danh ổn định và vòng đời thay đổi theo thời gian. Ví dụ gồm Node,
Model, Deployment, Command và Incident. Hai entity có cùng thuộc tính nhưng
khác định danh vẫn là hai đối tượng khác nhau.

\subsection{Value Object}

Value Object được xác định bởi giá trị và thường bất biến, chẳng hạn ModelVersion,
Checksum, DeviceCapability, MetricWindow hoặc DeploymentPolicy. Value Object
giúp đóng gói validation và tránh truyền các chuỗi rời rạc không có ngữ nghĩa.

\subsection{Aggregate}

Aggregate là biên nhất quán cho một nhóm entity và value object. Aggregate root
kiểm soát mọi thay đổi bên trong. Deployment có thể là aggregate quản lý target,
batch, trạng thái và policy; các thay đổi phải đi qua hành vi của Deployment thay
vì cập nhật bảng trực tiếp.

\subsection{Repository}

Repository cung cấp giao diện tải và lưu aggregate theo ngôn ngữ domain. Domain
không phụ thuộc chi tiết PostgreSQL, object storage hoặc message broker. Điều này
giúp kiểm thử quy tắc nghiệp vụ độc lập với hạ tầng.

\subsection{Domain Service}

Domain Service chứa hành vi nghiệp vụ không thuộc tự nhiên về một entity, ví dụ
chọn nhóm canary dựa trên capability hoặc đánh giá chính sách triển khai. Service
không nên trở thành nơi chứa toàn bộ logic; invariant vẫn được giữ trong aggregate
khi có thể.

\subsection{Bounded Context}

Bounded Context xác định biên mô hình và ngôn ngữ. Fleet quản lý node và
capability; Model Registry quản lý model/version/artifact; Deployment quản lý
rollout; Incident quản lý sự cố; Knowledge quản lý tài liệu; Copilot điều phối hội
thoại và tool call; Ops cung cấp quan sát vận hành. Các context trao đổi qua API
hoặc event với hợp đồng rõ ràng.
\section{Kết chương}

Chương 2 đã trình bày chuỗi kiến thức từ model đến platform. YOLO, pruning, QAT
và TensorRT cung cấp nền tảng cho nghiên cứu tối ưu; Model Registry, Deployment,
Edge Runtime và nền tảng phân tán giải thích cách đưa model vào vận hành; RAG,
AI Agent, MCP và DDD cung cấp cơ sở cho Knowledge, Copilot và cách chia domain.
Chương 3 sử dụng các cơ sở này để xây dựng và đánh giá pipeline tối ưu YOLO.